\BeleskeID{07-s005}
% element: 07-s005:d01
% element: 07-s005:e05
% element: 07-s005:e06
Za I matricu treba pitati koje direktne ili negirane ulazne promenljive mogu učestvovati u svakom proizvodu \(Y\), koliko različitih proizvoda \(m\) postoji i koliko ulaza ima svako I kolo. Za ILI matricu pitamo koji proizvodi mogu učestvovati u izlaznoj funkciji \(F\) i koliko proizvoda može ući u svako ILI kolo. Broj spoljašnjih ulaza je \(n\), a izlaza \(k\).

Pored toga da li se programira jedna ili obe matrice, bitno je koliko slobode postoji unutar svake matrice. Sam naziv programabilne komponente ne garantuje da se svaka veza može birati nezavisno.
